\chapter*{Bibliography}
\addcontentsline{toc}{chapter}{Bibliography}

\begin{thebibliography}{99}

\bibitem{gardner1970} Gardner, M. (1970).
  ``The fantastic combinations of John Conway's new solitaire game
  `life'.'' \emph{Scientific American} \textbf{223}(4), 120--123.

\bibitem{vonneumann1966} von Neumann, J. (1966).
  \emph{Theory of Self-Reproducing Automata} (edited and completed by
  A.~W.~Burks). University of Illinois Press.

\bibitem{shannon1948} Shannon, C.~E. (1948).
  ``A mathematical theory of communication.''
  \emph{Bell System Technical Journal} \textbf{27}, 379--423 and
  623--656.

\bibitem{wolfram1984} Wolfram, S. (1984).
  ``Universality and complexity in cellular automata.''
  \emph{Physica D} \textbf{10}, 1--35.

\bibitem{wolfram2002} Wolfram, S. (2002).
  \emph{A New Kind of Science}. Wolfram Media.

\bibitem{wuensche1992} Wuensche, A. and Lesser, M. (1992).
  \emph{The Global Dynamics of Cellular Automata: An Atlas of Basin of
  Attraction Fields of One-Dimensional Cellular Automata}.
  Addison-Wesley. (The atlas referred to in Section~\ref{sec:k5n21}.)

\bibitem{toffoli1987} Toffoli, T. and Margolus, N. (1987).
  \emph{Cellular Automata Machines: A New Environment for Modeling}.
  MIT Press.

\bibitem{langton1990} Langton, C.~G. (1990).
  ``Computation at the edge of chaos: phase transitions and emergent
  computation.'' \emph{Physica D} \textbf{42}, 12--37.

\bibitem{cook2004} Cook, M. (2004).
  ``Universality in elementary cellular automata.''
  \emph{Complex Systems} \textbf{15}, 1--40. (Proves Rule 110
  universal.)

\bibitem{berlekamp1982} Berlekamp, E.~R., Conway, J.~H., and Guy,
  R.~K. (1982). \emph{Winning Ways for Your Mathematical Plays},
  Vol.~2. Academic Press. (Universality of Life.)

\bibitem{lagarias2010} Lagarias, J.~C., ed. (2010).
  \emph{The Ultimate Challenge: The $3x+1$ Problem}.
  American Mathematical Society.

\bibitem{tao2019} Tao, T. (2022).
  ``Almost all orbits of the Collatz map attain almost bounded
  values.'' \emph{Forum of Mathematics, Pi} \textbf{10}, e12.
  (First circulated 2019.)

\bibitem{barina2020} Barina, D. (2021).
  ``Convergence verification of the Collatz problem.''
  \emph{The Journal of Supercomputing} \textbf{77}, 2681--2688.

\bibitem{sierpinski1915} Sierpi\'nski, W. (1915).
  ``Sur une courbe dont tout point est un point de ramification.''
  \emph{Comptes Rendus} \textbf{160}, 302--305.

\bibitem{harris2020} Harris, C.~R.\ et al. (2020).
  ``Array programming with NumPy.''
  \emph{Nature} \textbf{585}, 357--362.

\bibitem{virtanen2020} Virtanen, P.\ et al. (2020).
  ``SciPy 1.0: fundamental algorithms for scientific computing in
  Python.'' \emph{Nature Methods} \textbf{17}, 261--272.

\bibitem{hoffman2010} Hoffman, D.~D., Singh, M., and Mark, J. (2010).
  ``Natural selection and veridical perceptions.''
  \emph{Journal of Theoretical Biology} \textbf{266}(4), 529--539.

\bibitem{hoffman2019} Hoffman, D.~D. (2019).
  \emph{The Case Against Reality: Why Evolution Hid the Truth from Our
  Eyes}. W.~W.~Norton.

\bibitem{friedman2016} Friedman, D. and Sinervo, B. (2016).
  \emph{Evolutionary Games in Natural, Social, and Virtual Worlds}.
  Oxford University Press.

\bibitem{kahneman2011} Kahneman, D. (2011).
  \emph{Thinking, Fast and Slow}. Farrar, Straus and Giroux.

\bibitem{espericueta2020} Espericueta, R. (2020).
  ``Veridical Perception and Evolutionary Dynamics.'' Course study,
  Applied Machine Learning Lab, UC Santa Cruz.

\end{thebibliography}

\vspace{2em}
\noindent\textbf{Related links.}\quad
The complete Python source (\texttt{cadyn} package), the figure-
and appendix-generation scripts, and the LaTeX source of this book are
distributed together.  See the \texttt{README} for installation and a
guided tour.
